\documentclass[10pt,letterpaper]{article}

% ---------------------------------------------------------
% PAGE SIZE / MARGINS
% ---------------------------------------------------------
\usepackage[letterpaper,margin=1in]{geometry}

% ---------------------------------------------------------
% FONT: Times New Roman (serif) for text and headings
% ---------------------------------------------------------
\usepackage[T1]{fontenc}
\usepackage{mathptmx} % Times New Roman-compatible text and math font

% ---------------------------------------------------------
% PAGE NUMBERS
% ---------------------------------------------------------
\usepackage{fancyhdr}
\pagestyle{fancy}
\fancyhf{}
\fancyfoot[C]{\thepage}
\renewcommand{\headrulewidth}{0pt}

% ---------------------------------------------------------
% GRAPHICS: vector graphics preferred; PNG allowed (no JPG)
% Include figures as .pdf/.eps/.svg-converted-to-pdf (vector)
% or as 300+ dpi .png/.gif if a bitmap is unavoidable.
% ---------------------------------------------------------
\usepackage{graphicx}
\graphicspath{{./figures/}}

% ---------------------------------------------------------
% CAPTIONS: 8pt font for figure/table captions
% ---------------------------------------------------------
\usepackage[font=it]{caption}
\captionsetup{font={it,footnotesize}} % ~8pt equivalent (see explicit sizing below)
\DeclareCaptionFont{eightpt}{\fontsize{8}{9.5}\selectfont}
\captionsetup{font={eightpt,it}}

% ---------------------------------------------------------
% TABLES
% ---------------------------------------------------------
\usepackage{array}
\usepackage{longtable}
\usepackage{booktabs}
\renewcommand{\arraystretch}{1.3}

% ---------------------------------------------------------
% SECTION / TITLE FONT SIZES
% Title: 24pt | Body text: 10pt | Index terms: 9pt | Captions/refs: 8pt
% ---------------------------------------------------------
\usepackage{titlesec}
\titleformat{\section}{\normalfont\fontsize{14}{16}\bfseries}{\thesection}{1em}{}
\titleformat{\subsection}{\normalfont\fontsize{12}{14}\bfseries\itshape}{\thesubsection}{1em}{}
\titleformat{\subsubsection}{\normalfont\fontsize{11}{13}\bfseries}{\thesubsubsection}{1em}{}

% Custom command for 9pt index terms
\newcommand{\indexterms}[1]{{\fontsize{9}{11}\selectfont \textbf{Index Terms---} #1\par}}
\newcommand{\ucseparator}{\vspace{1em}\noindent\rule{\linewidth}{0.4pt}\vspace{1.5em}}

% Bibliography / references at 8pt
% Using standard thebibliography (no external bib backend required).
% If you prefer biblatex/biber, install biblatex and biber and replace
% this block with: \usepackage[backend=biber,style=ieee]{biblatex}
%                  \renewcommand*{\bibfont}{\fontsize{8}{9.5}\selectfont}
\makeatletter
\renewenvironment{thebibliography}[1]
 {\section*{References}%
  \fontsize{8}{9.5}\selectfont
  \list{\@biblabel{\@arabic\c@enumiv}}%
       {\settowidth\labelwidth{\@biblabel{#1}}%
        \leftmargin\labelwidth
        \advance\leftmargin\labelsep
        \@openbib@code
        \usecounter{enumiv}%
        \let\p@enumiv\@empty
        \renewcommand\theenumiv{\@arabic\c@enumiv}}%
  \sloppy
  \clubpenalty4000
  \@clubpenalty \clubpenalty
  \widowpenalty4000%
  \sfcode`\.\@m}
 {\def\@noitemerr
   {\@latex@warning{Empty `thebibliography' environment}}%
  \endlist}
\makeatother

\usepackage{hyperref}
\hypersetup{colorlinks=true, linkcolor=black, urlcolor=blue, citecolor=black}

\begin{document}

% ===========================================================
% TITLE (24pt)
% ===========================================================
\begin{center}
{\fontsize{24}{28}\selectfont \textbf{Healthy Living}}
\end{center}
\vspace{0.5em}

% Optional: Index terms line (9pt) -- uncomment and edit if needed
% \indexterms{use case model, actor-goal list, requirements}

% ===========================================================
% SECTION 1: ACTOR-GOAL LIST
% ===========================================================
\section*{1. Actor-Goal List}

\begin{table}[h!]
\centering
\caption{Actor-Goal List}
\label{tab:actor-goal-list}
\begin{tabular}{|p{0.35\linewidth}|p{0.5\linewidth}|}
\hline
\textbf{Actor} & \textbf{Goal} \\
\hline
Student (Primary End-User) & Manage and track personal daily intake through the web application \\
\hline
System Administrator (Admin) & Maintain, secure, and ensure the overall health and performance of the backend and web application \\
\hline
Content Moderator / Content Manager & Oversee the recipe database to ensure the quality, accuracy, and safety of platform content \\
\hline
Support Specialist & Handle customer service requests and resolve users' technical or functional issues \\
\hline
\end{tabular}
\end{table}

% ===========================================================
% SECTION 2: USE CASE MODEL
% ===========================================================
\section*{2. Use Case Model}

\subsection*{Id: UC-01}

\subsection*{Use Case: Create a Healthy Living Account}

\subsubsection*{Description}
A Student signs up for an account so they can start tracking their meals and health goals.

\subsubsection*{Level}
User-goal

\subsubsection*{Primary Actor}
Student

\subsubsection*{Supporting Actors}
None

\subsubsection*{Stakeholders and Interests}
\begin{itemize}
    \item \textbf{Student:} Wants an account so they can start using the app.
    \item \textbf{System Administrator:} Wants accounts created safely and stored correctly.
    \item \textbf{Support Specialist:} Wants sign-up to go smoothly so people don't need help.
\end{itemize}

\subsubsection*{Pre-Conditions}
The Student doesn't already have an account.

\subsubsection*{Post Conditions}

\textbf{Success end condition}\\
The Student has a new account and can start using the app.

\textbf{Failure end condition}\\
No account is created.

\subsection*{Main Success Scenario}
\begin{enumerate}
    \item The Student asks to sign up.
    \item The System asks for information about the Student.
    \item The Student enters their information.
    \item The System checks the information is valid and not already in use.
    \item The Student agrees to the terms.
    \item The System creates the account and informs the Student
\end{enumerate}

\subsection*{Extensions}
\begin{enumerate}
    \item[3a.] Student gives up before finishing.
    \begin{enumerate}
        \item[1.] Use case ends, no account created.
    \end{enumerate}
    \item[4a.] Info is missing or wrong.
    \begin{enumerate}
        \item[1.] System informs the Student what is wrong.
        \item[2.] Student fixes it and use case continues at step 4.
    \end{enumerate}
    \item[4b.] Account already exists for that Student.
    \begin{enumerate}
        \item[1.] System lets them know and suggests logging in instead.
        \item[2.] Use case ends.
    \end{enumerate}
    \item[5a.] Student doesn't agree to the terms.
    \begin{enumerate}
        \item[1.] Use case ends, no account created.
    \end{enumerate}
\end{enumerate}

\subsection*{Special Requirements}
See Supplementary Specification for Usability, Reliability, Performance and Supportability requirements.

\ucseparator

\subsection*{Id: UC-02}

\subsection*{Use Case: Maintain Personal Profile}

\subsubsection*{Description}
A Student updates their profile information, so their goals and preferences stay up to date.

\subsubsection*{Level}
User-goal

\subsubsection*{Primary Actor}
Student

\subsubsection*{Supporting Actors}
None

\subsubsection*{Stakeholders and Interests}
\begin{itemize}
    \item \textbf{Student:} Wants their profile to reflect their current goals and preferences.
    \item \textbf{Content Moderator:} Wants profile info to stay consistent with platform content and guidelines.
\end{itemize}

\subsubsection*{Pre-Conditions}
The Student is logged in and has an existing account.

\subsubsection*{Post Conditions}

\textbf{Success end condition}\\
The Student's profile reflects the changes they made.

\textbf{Failure end condition}\\
The profile stays as it was before.

\subsection*{Main Success Scenario}
\begin{enumerate}
    \item The Student asks to update their profile.
    \item The System shows the current information of the Student.
    \item The Student makes the changes they want.
    \item The Student submits the changes.
    \item The System checks if the changes are valid.
    \item The System saves the changes and informs the Student.
\end{enumerate}

\subsection*{Extensions}
\begin{enumerate}
    \item[3a.] Student decides not to make any changes.
    \begin{enumerate}
        \item[1.] Use case ends, profile stays the same.
    \end{enumerate}
    \item[5a.] Some of the changes aren't valid.
    \begin{enumerate}
        \item[1.] System tells the Student what's wrong.
        \item[2.] Student fixes it and use case continues at step 4.
    \end{enumerate}
\end{enumerate}

\subsection*{Special Requirements}
See Supplementary Specification for Usability, Reliability, Performance and Supportability requirements.

\ucseparator

\subsection*{Id: UC-03}
\subsection*{Use Case: Plan Weekly Meals}
\subsubsection*{Description}
A Student creates a meal plan for the week so they know what they'll eat and can stay on track with their goals.
\subsubsection*{Level}
User-goal
\subsubsection*{Primary Actor}
Student
\subsubsection*{Supporting Actors}
None
\subsubsection*{Stakeholders and Interests}
\begin{itemize}
    \item \textbf{Student:} Wants a plan for the week to achieve his goals and preferences.
    \item \textbf{Content Moderator:} Wants recipes suggested through the plan to be accurate.
\end{itemize}

\subsubsection*{Pre-Conditions}
The Student is logged in and has a profile with their goals and preferences set up.

\subsubsection*{Post Conditions}
\textbf{Success end condition}\\
The Student has a weekly meal plan saved to their account.

\textbf{Failure end condition}\\
No new meal plan is saved.

\subsection*{Main Success Scenario}
\begin{enumerate}
    \item The Student asks to plan meals for the week.
    \item The System suggests recipes based on the Student's goals and preferences.
    \item The Student picks meals for each day.
    \item The Student reviews the plan.
    \item The Student confirms the plan.
    \item The System saves the weekly meal plan and confirms it worked.
\end{enumerate}
\subsection*{Extensions}
\begin{enumerate}
    \item[2a.] System has no suitable recipes to suggest.
    \begin{enumerate}
        \item[1.] System lets the Student know and asks them to adjust their preferences or search for recipes directly.
        \item[2.] Student updates their preferences or searches for and picks recipes themself.
        \item[3.] Use case continues at step 3.
    \end{enumerate}
    \item[3a.] Student doesn't like any of the suggested meals or wants more recipes to choose from.
    \begin{enumerate}
        \item[1.] Student searches for other recipes and picks them for the relevant days.
        \item[2.] Use case continues at step 3.
    \end{enumerate}
    \item[4a.] Student leaves the plan incomplete (some days without meals).
    \begin{enumerate}
        \item[1.] Student decides to finish later.
        \item[2.] Use case ends, no plan is saved.
    \end{enumerate}
    \item[5a.] Student decides not to save the plan.
    \begin{enumerate}
        \item[1.] Use case ends, no plan is saved.
    \end{enumerate}
\end{enumerate}
\subsection*{Special Requirements}
See Supplementary Specification for Usability, Reliability, Performance and Supportability requirements.

\ucseparator

\subsection*{Id: UC-04}

\subsection*{Use Case: Manage Grocery Inventory}

\subsubsection*{Description}
A Student keeps track of what groceries they have, adding items as they buy them and removing items as they use them.

\subsubsection*{Level}
User-goal

\subsubsection*{Primary Actor}
Student

\subsubsection*{Supporting Actors}
None

\subsubsection*{Stakeholders and Interests}
\begin{itemize}
    \item \textbf{Student:} Wants an accurate picture of what groceries they have so they can plan meals.
\end{itemize}

\subsubsection*{Pre-Conditions}
The Student is logged in and has an existing account.

\subsubsection*{Post Conditions}

\textbf{Success end condition}\\
The grocery inventory reflects the changes they made.

\textbf{Failure end condition}\\
The inventory data doesn't change.

\subsection*{Main Success Scenario}
\begin{enumerate}
    \item The Student asks to update their grocery inventory.
    \item The System shows the Student's current inventory.
    \item The Student adds new items they've bought, or marks items as used.
    \item The Student submits the changes.
    \item The System checks the changes are valid.
    \item The System saves the updated inventory and informs the Student.
\end{enumerate}

\subsection*{Extensions}
\begin{enumerate}
    \item[3a.] Student decides not to make any changes.
    \begin{enumerate}
        \item[1.] Use case ends, inventory stays the same.
    \end{enumerate}
    \item[3b.] Student tries to remove more of an item than they have on hand.
    \begin{enumerate}
        \item[1.] System lets the Student know and adjusts the amount to zero.
        \item[2.] Use case continues at step 4.
    \end{enumerate}
    \item[3c.] An item marked as used is needed for a meal in the Student's current meal plan.
    \begin{enumerate}
        \item[1.] System notifies the Student that a planned meal is affected.
        \item[2.] Use case continues at step 4, or at UC08 if the Student chooses to review the plan.
    \end{enumerate}
    \item[5a.] Some of the changes aren't valid.
    \begin{enumerate}
        \item[1.] System tells the Student what's wrong.
        \item[2.] Student fixes it and use case continues at step 4.
    \end{enumerate}
\end{enumerate}

\subsection*{Special Requirements}
See Supplementary Specification for Usability, Reliability, Performance and Supportability requirements.

\ucseparator

\subsection*{Id: UC-05}

\subsection*{Use Case: Find Recipes}

\subsubsection*{Description}
A Student searches for recipes that match what they're looking for, based on ingredients or dietary preferences.

\subsubsection*{Level}
User-goal

\subsubsection*{Primary Actor}
Student

\subsubsection*{Supporting Actors}
None

\subsubsection*{Stakeholders and Interests}
\begin{itemize}
    \item \textbf{Student:} Wants to find recipes that fit their needs and preferences.
    \item \textbf{Content Moderator:} Wants search results to only include accurate and approved recipes.
\end{itemize}

\subsubsection*{Pre-Conditions}
The Student is logged in.

\subsubsection*{Post Conditions}

\textbf{Success end condition}\\
The Student sees a list of recipes matching their search.

\textbf{Failure end condition}\\
No search results are shown.

\subsection*{Main Success Scenario}
\begin{enumerate}
    \item The Student asks to search for recipes.
    \item The Student enters their search criteria (e.g., ingredients and dietary needs).
    \item The System looks for recipes that match the criteria.
    \item The System displays a list of matching recipes.
    \item The Student selects a recipe to view.
    \item The System shows the recipe details.
\end{enumerate}

\subsection*{Extensions}
\begin{enumerate}
    \item[3a.] No recipes match the criteria.
    \begin{enumerate}
        \item[1.] System informs the Student and suggests adjusting the search.
        \item[2.] Student changes their criteria and use case continues at step 2.
    \end{enumerate}
    \item[5a.] Student doesn't select a recipe.
    \begin{enumerate}
        \item[1.] Use case ends, no recipe is viewed.
    \end{enumerate}
\end{enumerate}

\subsection*{Special Requirements}
See Supplementary Specification for Usability, Reliability, Performance and Supportability requirements.

\ucseparator

\subsection*{Id: UC-06}

\subsection*{Use Case: Obtain Recipe Recommendations}

\subsubsection*{Description}
A Student asks the system for recipe suggestions based on their goals, preferences and ingredients he has.

\subsubsection*{Level}
User-goal

\subsubsection*{Primary Actor}
Student

\subsubsection*{Supporting Actors}
None

\subsubsection*{Stakeholders and Interests}
\begin{itemize}
    \item \textbf{Student:} Wants relevant recipe ideas without having to search for them manually.
    \item \textbf{Content Moderator:} Wants recommended recipes to be accurate and approved.
\end{itemize}

\subsubsection*{Pre-Conditions}
The Student is logged in and has a profile with their goals and preferences set up.

\subsubsection*{Post Conditions}

\textbf{Success end condition}\\
The Student sees a list of recommended recipes.

\textbf{Failure end condition}\\
No recommendations are shown.

\subsection*{Main Success Scenario}
\begin{enumerate}
    \item The Student asks for recipe recommendations.
    \item The System looks at the Student's goals, preferences and inventory.
    \item The System finds recipes that fit.
    \item The System shows the Student the recommended recipes.
    \item The Student selects a recipe to view.
    \item The System shows the recipe details.
\end{enumerate}

\subsection*{Extensions}
\begin{enumerate}
    \item[3a.] System can't find any recipes that fit.
    \begin{enumerate}
        \item[1.] System lets the Student know and suggests updating their preferences.
        \item[2.] Use case ends, no recommendations shown.
    \end{enumerate}
    \item[5a.] Student doesn't select a recipe.
    \begin{enumerate}
        \item[1.] Use case ends, no recipe is viewed.
    \end{enumerate}
\end{enumerate}

\subsection*{Special Requirements}
See Supplementary Specification for Usability, Reliability, Performance and Supportability requirements.

\ucseparator

\subsection*{Id: UC-07}

\subsection*{Use Case: Generate Shopping List}

\subsubsection*{Description}
A Student gets a shopping list based on their meal plan and what they already have in their inventory.

\subsubsection*{Level}
User-goal

\subsubsection*{Primary Actor}
Student

\subsubsection*{Supporting Actors}
None

\subsubsection*{Stakeholders and Interests}
\begin{itemize}
    \item \textbf{Student:} Wants an accurate list of what to buy without having to figure it out manually.
    \item \textbf{System Administrator:} Wants shopping list data stored reliably.
\end{itemize}

\subsubsection*{Pre-Conditions}
The Student is logged in and has a saved weekly meal plan.

\subsubsection*{Post Conditions}

\textbf{Success end condition}\\
The Student has a shopping list saved to their account.

\textbf{Failure end condition}\\
No shopping list is saved.

\subsection*{Main Success Scenario}
\begin{enumerate}
    \item The Student asks to generate a shopping list.
    \item The System looks at the ingredients needed for the meal plan.
    \item The System compares them against the Student's current inventory.
    \item The System puts together a list of the items still needed.
    \item The System shows the Student the shopping list.
    \item The Student reviews and confirms the list.
    \item The System saves the shopping list and confirms it worked.
\end{enumerate}

\subsection*{Extensions}
\begin{enumerate}
    \item[4a.] Nothing is needed because the Student already has everything.
    \begin{enumerate}
        \item[1.] System lets the Student know no shopping is needed.
        \item[2.] Use case ends, no list is saved.
    \end{enumerate}
    \item[6a.] Student wants to add or remove items from the list.
    \begin{enumerate}
        \item[1.] Student makes changes to the list.
        \item[2.] Use case continues at step 6.
    \end{enumerate}
    \item[6b.] Student decides not to save the list.
    \begin{enumerate}
        \item[1.] Use case ends, no list is saved.
    \end{enumerate}
\end{enumerate}

\subsection*{Special Requirements}
See Supplementary Specification for Usability, Reliability, Performance and Supportability requirements.

\ucseparator

\subsection*{Id: UC-08}

\subsection*{Use Case: Review Meal Plan}

\subsubsection*{Description}
A Student looks over their current weekly meal plan to check it still fits their needs and makes adjustments if it doesn't.

\subsubsection*{Level}
User-goal

\subsubsection*{Primary Actor}
Student

\subsubsection*{Supporting Actors}
None

\subsubsection*{Stakeholders and Interests}
\begin{itemize}
    \item \textbf{Student:} Wants to make sure their plan still matches their goals and make changes if it doesn't.
    \item \textbf{Content Moderator:} Wants recipes to be accurate and approved.
\end{itemize}

\subsubsection*{Pre-Conditions}
The Student is logged in and has a saved weekly meal plan.

\subsubsection*{Post Conditions}

\textbf{Success end condition}\\
The meal plan reflects any changes the Student made or is confirmed as-is.

\textbf{Failure end condition}\\
The meal plan stays as it was before.

\subsection*{Main Success Scenario}
\begin{enumerate}
    \item The Student asks to review their meal plan.
    \item The System shows the Student their current weekly plan.
    \item The Student checks the plan against their goals.
    \item The Student makes changes to one or more meals.
    \item The Student confirms the updated plan.
    \item The System saves the changes and informs the Student.
\end{enumerate}

\subsection*{Extensions}
\begin{enumerate}
    \item[3a.] Student decides the plan is fine as-is.
    \begin{enumerate}
        \item[1.] Use case ends, no changes made.
    \end{enumerate}
    \item[4a.] Student wants to swap out a meal.
    \begin{enumerate}
        \item[1.] Student picks a different recipe for that meal.
        \item[2.] Use case continues at step 4.
    \end{enumerate}
    \item[5a.] Student decides not to keep the changes.
    \begin{enumerate}
        \item[1.] Use case ends, plan stays the same.
    \end{enumerate}
\end{enumerate}

\subsection*{Special Requirements}
See Supplementary Specification for Usability, Reliability, Performance and Supportability requirements.


\ucseparator

\subsection*{Id: UC-09}

\subsection*{Use Case: Moderate Recipe Content}

\subsubsection*{Description}
A Content Moderator reviews recipes submitted or flagged on the platform to make sure they're accurate and safe before they stay live.

\subsubsection*{Level}
User-goal

\subsubsection*{Primary Actor}
Content Moderator

\subsubsection*{Supporting Actors}
None

\subsubsection*{Stakeholders and Interests}
\begin{itemize}
    \item \textbf{Content Moderator:} Wants an efficient way to review and act on recipes needing attention.
    \item \textbf{Student:} Wants the recipes they see on the platform to be accurate and safe.
\end{itemize}

\subsubsection*{Pre-Conditions}
The Content Moderator is logged in and there are recipes pending review.

\subsubsection*{Post Conditions}

\textbf{Success end condition}\\
The reviewed recipe is approved, edited or removed based on the Moderator's decision.

\textbf{Failure end condition}\\
The recipe's status stays unchanged.

\subsection*{Main Success Scenario}
\begin{enumerate}
    \item The Content Moderator asks to review pending recipes.
    \item The System shows a list of recipes needing review.
    \item The Content Moderator selects a recipe to review.
    \item The System shows the recipe's details.
    \item The Content Moderator checks the recipe for accuracy and safety.
    \item The Content Moderator approves the recipe.
    \item The System updates the recipe's status and confirms it worked.
\end{enumerate}

\subsection*{Extensions}
\begin{enumerate}
    \item[5a.] Content Moderator finds issues that can be fixed.
    \begin{enumerate}
        \item[1.] Content Moderator edits the recipe.
        \item[2.] Use case continues at step 6.
    \end{enumerate}
    \item[5b.] Content Moderator finds the recipe unsafe or inappropriate.
    \begin{enumerate}
        \item[1.] Content Moderator removes the recipe.
        \item[2.] Use case continues at step 7.
    \end{enumerate}
    \item[6a.] Content Moderator decides to leave the recipe as pending for further review.
    \begin{enumerate}
        \item[1.] Use case ends, recipe stays pending.
    \end{enumerate}
\end{enumerate}

\subsection*{Special Requirements}
See Supplementary Specification for Usability, Reliability, Performance and Supportability requirements.

\ucseparator

\subsection*{Id: UC-10}

\subsection*{Use Case: Resolve User Support Requests}

\subsubsection*{Description}
A Support Specialist addresses a user's support request to help them resolve a technical or functional issue.

\subsubsection*{Level}
User-goal

\subsubsection*{Primary Actor}
Support Specialist

\subsubsection*{Supporting Actors}
Student

\subsubsection*{Stakeholders and Interests}
\begin{itemize}
    \item \textbf{Support Specialist:} Wants to understand and resolve Student issues.
    \item \textbf{Student:} Wants their issue resolved quickly and correctly.
\end{itemize}

\subsubsection*{Pre-Conditions}
The Support Specialist is logged in and there is an open support request.

\subsubsection*{Post Conditions}

\textbf{Success end condition}\\
The support request is resolved and closed.

\textbf{Failure end condition}\\
The support request stays open.

\subsection*{Main Success Scenario}
\begin{enumerate}
    \item The Support Specialist asks to view open support requests.
    \item The System shows a list of open requests.
    \item The Support Specialist selects a request to work on.
    \item The System shows the request details.
    \item The Support Specialist investigates the issue.
    \item The Support Specialist resolves the issue.
    \item The Support Specialist closes the request.
    \item The System updates the request status and informs the Student.
\end{enumerate}

\subsection*{Extensions}
\begin{enumerate}
    \item[5a.] Support Specialist needs more information from the Student.
    \begin{enumerate}
        \item[1.] Support Specialist contacts the Student for details.
        \item[2.] Student provides the requested details.
        \item[3.] Use case continues at step 5.
    \end{enumerate}
    \item[6a.] Support Specialist can't resolve the issue directly.
    \begin{enumerate}
        \item[1.] Support Specialist escalates the request.
        \item[2.] Use case ends, request stays open pending escalation.
    \end{enumerate}
\end{enumerate}

\subsection*{Special Requirements}
See Supplementary Specification for Usability, Reliability, Performance and Supportability requirements.

\ucseparator

\subsection*{Id: UC-11}

\subsection*{Use Case: Manage User Accounts}

\subsubsection*{Description}
A System Administrator makes changes to user accounts, such as suspending, reinstating or removing them, to keep the platform secure and running properly.

\subsubsection*{Level}
User-goal

\subsubsection*{Primary Actor}
System Administrator

\subsubsection*{Supporting Actors}
None

\subsubsection*{Stakeholders and Interests}
\begin{itemize}
    \item \textbf{System Administrator:} Wants to manage accounts and keep the platform secure.
    \item \textbf{Student:} Wants their account managed securely.
\end{itemize}

\subsubsection*{Pre-Conditions}
The System Administrator is logged in and the account to manage exists.

\subsubsection*{Post Conditions}

\textbf{Success end condition}\\
The account reflects the change the Administrator made.

\textbf{Failure end condition}\\
The account stays unchanged.

\subsection*{Main Success Scenario}
\begin{enumerate}
    \item The System Administrator asks to view user accounts.
    \item The System shows a list of accounts.
    \item The System Administrator selects an account to manage.
    \item The System shows the account's details.
    \item The System Administrator chooses an action to take on the account.
    \item The System Administrator confirms the action.
    \item The System applies the change and informs the system administrator.
\end{enumerate}

\subsection*{Extensions}
\begin{enumerate}
    \item[5a.] Administrator decides not to take any action.
    \begin{enumerate}
        \item[1.] Use case ends, account stays unchanged.
    \end{enumerate}
    \item[6a.] Administrator cancels before confirming.
    \begin{enumerate}
        \item[1.] Use case ends, account stays unchanged.
    \end{enumerate}
\end{enumerate}

\subsection*{Special Requirements}
See Supplementary Specification for data security and account privacy requirements.
\clearpage

% ===========================================================
% SECTION 3: USE CASE DIAGRAM
% ===========================================================

\section*{3. Use Case Diagram}

\begin{figure}[h!]
\centering
\includegraphics[width=0.9\linewidth]{images/UseCaseDiagram.png}
\caption{Use case diagram for the Healthy Living App.}
\label{fig:use-case-diagram}
\end{figure}
\clearpage

% ===========================================================
% SECTION 4: SEQUENCE DIAGRAM - UC04
% ===========================================================
\section*{4. Sequence Diagram: Manage Grocery Inventory (UC-04)}
 
\begin{figure}[h!]
\centering
\includegraphics[width=0.9\linewidth]{images/SequenceDiagramUC04.png}
\caption{Sequence diagram for the Manage Grocery Inventory use case (UC-04).}
\label{fig:sequence-diagram-uc04}
\end{figure}
\clearpage

% ===========================================================
% SECTION 5: ACTIVITY DIAGRAM - UC03
% ===========================================================
\section*{5. Activity Diagram: Plan Weekly Meals (UC-03)}
 
\begin{figure}[h!]
\centering
\includegraphics[width=0.9\linewidth]{images/ActivityDiagramUC03.png}
\caption{Activity diagram for the Plan Weekly Meals use case (UC-03).}
\label{fig:activity-diagram-uc03}
\end{figure}
 
\clearpage

% ===========================================================
% EXAMPLE FIGURE (vector graphic, 300+ dpi if bitmap, captioned)
% ===========================================================
% \begin{figure}[h!]
% \centering
% \includegraphics[width=0.8\linewidth]{diagram.pdf} % vector format preferred
% \caption{Use case diagram showing actor interactions.}
% \label{fig:use-case-diagram}
% \end{figure}

% ===========================================================
% REFERENCES (8pt, automatically applied by the thebibliography
% environment redefinition above)
% ===========================================================
% \begin{thebibliography}{9}
% \bibitem{example2024} A. Author, ``Title of the paper,'' \textit{Journal Name}, vol. 1, no. 1, pp. 1--10, 2024.
% \end{thebibliography}

\end{document}
